% Algoritmi

Алгоритми за рад са динамичким графовима имају кључну улогу у ефикасном одржавању и обради информација у графовима чија се структура мења током времена. За разлику од алгоритама над статичким графовима, алгоритми за динамичке графове често, поред саме алгоритамске процедуре, захтевају и употребу специјализованих структура података које омогућавају ефикасно ажурирање и одржавање информација од значаја. Због тога су овакви алгоритми по правилу сложенији и концептуално другачији у односу на класичне алгоритме над статичким графовима који се најчешће срећу у литератури.

У овом поглављу разматрају се проблеми одржавања информација о повезаности чворова и одржавања минималног разапињућег стабла у динамичким графовима. Ови проблеми имају значајну улогу у теорији графова и бројне примене у различитим областима, укључујући рачунарске мреже, друштвене мреже и биоинформатику.

\section{Динамичка повезаност}

Одржавање информације о повезаности чворова један је од најпроучаванијих проблема у области динамичких графова. Основни упит који је потребан за решавање овог проблема је \texttt{isConnected(u, v)}, који проверава да ли постоји пут између чворова \texttt{u} и \texttt{v} \cite{APPLICATIONS_OF_DYNAMIC_GRAPH}. У литератури се могу пронаћи ефикасне структуре података које омогућавају одржавање информација о повезаности у динамичким графовима, као што су \textit{Link-Cut Tree}, \textit{Euler Tour Tree} и \textit{Top Tree} \cite{APPLICATIONS_OF_DYNAMIC_GRAPH, holmPolylogarithmicDeterministicFullydynamic2001}. Ове структуре података омогућавају веома ефикасно одржавање информација о повезаности у динамичким графовима, али су по правилу сложене за имплементацију и најчешће специјализоване за одређену групу упита и операција над графом.

У овом поглављу биће предстваљено проширење DHB структуре података која омогућава одржавање информација о повезаности у инкременталним неусмереним динамичким графовима. Ово проширење засновано је на дисјунктним скуповима и омогућава ефикасно одржавање информација о компонентама повезаности у динамичким графовима. Упит \texttt{isConnected(u, v)} може се одговорити у константном времену. Додавање чворова има константу сложеност јер дирекнот користи DHB, док је додавање грана $O(\log n)$ због потребе за спајањем компоненти повезаности. Ово проширење DHB структуре података представља ефикасно решење за одржавање информација о повезаности у инкременталним неусмереним динамичким графовима.

\subsection{Имплементација}

\begin{lstlisting}[style=thesiscpp, caption={Интерфејс класе \texttt{DynamicConnectivityIncremental}}, label={lst:dc_interface}]
class DynamicConnectivityIncremental {
private:
    dg::Graph graph;
    std::vector<int> parent;
    std::vector<int> rank;

    int find(int x);
    void unite(int a, int b);

public:
    DynamicConnectivityIncremental() = default;

    int VertexCount() const;
    dg::VertexID AddVertex();
    
    void AddEdge(dg::VertexID vertex_id1, dg::VertexID vertex_id2, dg::Weight weight = 1);
    dg::Weight GetEdgeWeight(dg::VertexID vertex_id1, dg::VertexID vertex_id2) const;
    void UpdateEdgeWeight(dg::VertexID vertex_id1, dg::VertexID vertex_id2, dg::Weight new_weight);

    bool IsConnected(dg::VertexID vertex_id1, dg::VertexID vertex_id2);
};
\end{lstlisting}

\section{Динамичко минимално разапињуће стабло}

% Algoritmi